% =====================================================
% 题型：半角公式与余弦定理选择题 (q_tjw_20260606_20_half_angle_cosine_rule.tex)
% =====================================================
% =====================================================
% 难度：★★ (中档)
% =====================================================
\newcommand{\qTriangleHalfAngleCosineRuleStem}{%
  在 \(\triangle ABC\) 中，\(\cos \dfrac{C}{2} = \dfrac{\sqrt{5}}{5}\)，\(BC = 1\)，\(AC = 5\)，则 \(AB =\) \hfill （\quad）%
}

\newcommand{\qTriangleHalfAngleCosineRuleOptions}{%
  \mcoptions[2]
    {\(4\sqrt{2}\)}
    {\(\sqrt{30}\)}
    {\(\sqrt{29}\)}
    {\(2\sqrt{5}\)}%
}

\newcommand{\qTriangleHalfAngleCosineRuleAnswer}{A}

\newcommand{\qTriangleHalfAngleCosineRuleSolution}{%
  因为 \(C \in (0,\pi)\)，所以 \(\dfrac{C}{2} \in \left(0,\dfrac{\pi}{2}\right)\)。
  由半角公式
  \[
    \cos C = 2\cos^2\frac{C}{2} - 1
    = 2 \cdot \frac{1}{5} - 1
    = -\frac{3}{5}.
  \]

  由余弦定理得
  \[
    AB^2
    = AC^2 + BC^2 - 2 \cdot AC \cdot BC \cos C
    = 25 + 1 - 2 \cdot 5 \cdot 1 \cdot \left(-\frac{3}{5}\right)
    = 32.
  \]
  所以 \(AB = 4\sqrt{2}\)。
  故选 A。%
}

\newcommand{\qTriangleHalfAngleCosineRuleDiagnosis}{%
  （留白待收集学生作答后补充）%
}

\newcommand{\qTriangleHalfAngleCosineRuleTeachingNotes}{%
  精讲候选。半角给的是 \(\cos \frac{C}{2}\)，必须先还原 \(\cos C\)，再进入余弦定理。课堂可追问：为什么不能把 \(\frac{C}{2}\) 直接当作夹角代入？%
}


% =====================================================
% 别名兼容：旧课次宏定义
% =====================================================
\newcommand{\qTJWZeroSixTwentyStem}{\qTriangleHalfAngleCosineRuleStem}
\newcommand{\qTJWZeroSixTwentyOptions}{\qTriangleHalfAngleCosineRuleOptions}
\newcommand{\qTJWZeroSixTwentyAnswer}{\qTriangleHalfAngleCosineRuleAnswer}
\newcommand{\qTJWZeroSixTwentySolution}{\qTriangleHalfAngleCosineRuleSolution}
\newcommand{\qTJWZeroSixTwentyDiagnosis}{\qTriangleHalfAngleCosineRuleDiagnosis}
\newcommand{\qTJWZeroSixTwentyTeachingNotes}{\qTriangleHalfAngleCosineRuleTeachingNotes}
